\documentclass[12pt,a4paper]{article}

\usepackage[italian]{babel}
\usepackage[T1]{fontenc}
\usepackage[utf8]{inputenc}
\usepackage{geometry}
\usepackage{microtype}
\usepackage{setspace}
\usepackage{parskip}

\geometry{margin=2.4cm}
\setstretch{1.08}

\title{\textbf{Riassunto della tesi}\\
\large Sviluppo di un sistema RAG aziendale e valutazione di LLM locali}
\author{Francesco Romeo}
\date{Anno Accademico 2025/2026}

\begin{document}

\maketitle

\section*{Contesto e obiettivo del lavoro}

La tesi affronta la progettazione, l'implementazione e la valutazione di un
sistema di \textit{Retrieval-Augmented Generation} (RAG) integrato in i3Guild,
la piattaforma usata da Intré S.r.l. per gestire le Gilde aziendali. Il punto
di partenza è un problema concreto: nel corso degli anni i3Guild ha raccolto
proposte, descrizioni, obiettivi, risultati, partecipanti e metadati relativi
alle attività di apprendimento interno dell'azienda. Questo patrimonio
informativo rappresenta una memoria organizzativa utile, ma non è sempre facile
da consultare con strumenti tradizionali, perché le informazioni sono
distribuite tra database relazionali, campi testuali, metadati e contenuti
collegati.

L'obiettivo del lavoro non è addestrare un nuovo modello linguistico. La tesi
studia invece come integrare componenti esistenti in un prototipo enterprise
on-premise, capace di rispondere a domande in linguaggio naturale sul corpus
delle Gilde. La soluzione deve rispettare tre vincoli principali. Il primo è la
riservatezza: i dati aziendali devono rimanere entro un perimetro controllato.
Il secondo è la sostenibilità infrastrutturale: il sistema deve poter funzionare
su risorse hardware realistiche, senza assumere la disponibilità di server GPU
dedicati. Il terzo è la manutenibilità: l'integrazione deve inserirsi
nell'applicazione esistente senza confondere sorgente dati, indicizzazione,
retrieval e generazione.

Il paradigma RAG è adatto a questo scenario perché separa la conoscenza usata
dal sistema dai parametri del modello generativo. Un Large Language Model
pre-addestrato può produrre testo coerente, ma non conosce i dati privati e
aggiornati dell'organizzazione. In un sistema RAG, invece, la domanda
dell'utente viene prima usata per recuperare documenti pertinenti da una base di
conoscenza esterna; solo dopo questi documenti vengono inseriti nel prompt del
modello, che genera una risposta fondata sul contesto recuperato. Questa
impostazione riduce il rischio di risposte scollegate dal corpus e permette di
aggiornare la conoscenza intervenendo sull'indice, non sui pesi del modello.

\section*{Fondamenti tecnici e requisiti}

La prima parte della tesi introduce i fondamenti necessari per comprendere il
prototipo. I Transformer costituiscono la base degli attuali modelli
linguistici e dei modelli di embedding. Il meccanismo di self-attention consente
di costruire rappresentazioni contestuali del testo, superando diversi limiti
delle architetture ricorrenti. Tuttavia, gli LLM rimangono sistemi a conoscenza
chiusa quando sono usati senza strumenti esterni: non hanno accesso nativo ai
dati interni, possono generare informazioni plausibili ma errate e non offrono
da soli un meccanismo di verifica sulle fonti aziendali.

Da questi limiti deriva la scelta del RAG. La tesi distingue due pipeline: una
pipeline di ingestione, eseguita offline, e una pipeline di inferenza, attivata
in tempo reale dalle richieste utente. La prima normalizza i contenuti,
costruisce documenti adatti al retrieval, genera embedding densi e sparsi e li
salva in un database vettoriale. La seconda riceve la query, produce le
rappresentazioni della domanda, recupera i documenti più pertinenti e costruisce
il prompt aumentato per il modello generativo. In questo modo il sistema può
combinare similarità semantica, segnali lessicali e filtri basati sui metadati.

Il caso d'uso scelto riguarda il supporto alla proposta di nuove Gilde. Prima
di proporre un nuovo tema, una persona può avere bisogno di sapere se argomenti
simili siano già stati affrontati, quali risultati siano stati prodotti, chi
abbia partecipato e in quale epoca siano nate iniziative affini. L'analisi del
corpus mostra che la Gilda è l'unità informativa più adatta
all'indicizzazione: contiene campi descrittivi, obiettivi, rischi, risultati
attesi o raggiunti e collegamenti con partecipanti ed epoch. I metadati non
vengono trattati come semplice testo, ma restano disponibili per filtrare e
contestualizzare le risposte.

I requisiti del sistema riflettono questa struttura. Il prototipo deve
recuperare informazioni dal corpus i3Guild, generare risposte in linguaggio
naturale, indicare il contesto usato, supportare filtri strutturati e integrarsi
con l'interfaccia Next.js esistente. Deve anche evitare dipendenze non
necessarie da servizi esterni, perché il caso di studio richiede controllo sui
dati e sull'infrastruttura. La scelta tecnologica cade quindi su PostgreSQL come
sorgente autoritativa, Qdrant come database vettoriale, Haystack come framework
di orchestrazione, FastAPI come microservizio RAG, Ollama per l'inferenza
locale e Next.js per l'integrazione applicativa.

\section*{Progettazione e implementazione del sistema}

L'architettura proposta separa in modo netto i componenti. PostgreSQL conserva i
dati applicativi di i3Guild; una pipeline di embedding legge le informazioni
rilevanti, normalizza i campi testuali e costruisce documenti Haystack; Qdrant
memorizza embedding densi, rappresentazioni sparse e payload strutturati; il
microservizio FastAPI espone retrieval, chat sincrona e chat in streaming;
l'applicazione Next.js consuma il servizio tramite client dedicati. Questa
suddivisione riduce l'accoppiamento: l'applicazione rimane la sorgente
autoritativa, l'indice può essere ricostruito o aggiornato, e il modello
generativo può essere sostituito senza riscrivere l'intero sistema.

La pipeline di ingestione è progettata come job batch. Prima controlla la
configurazione, la connessione a PostgreSQL, la raggiungibilità di Qdrant e la
coerenza dei parametri di embedding. Poi estrae le Gilde con le relative epoch e
i partecipanti, pulisce i campi HTML con una normalizzazione testuale e compone
un documento strutturato in sezioni riconoscibili. In condizioni ordinarie una
Gilda viene indicizzata come documento consolidato, perché la dimensione media
dei campi resta compatibile con la finestra del modello di embedding. Per gli
outlier è prevista una soglia di sicurezza: se il testo supera il limite, viene
suddiviso in chunk con sovrapposizione controllata. Gli identificativi dei
documenti sono deterministici, così una reindicizzazione non duplica i punti in
Qdrant ma sovrascrive quelli esistenti.

Il retrieval non si limita a una ricerca vettoriale. La query dell'utente viene
normalizzata e, quando possibile, riscritta per isolare il tema effettivo della
domanda. Nel dominio i3Guild parole come "gilda", "fondare" o "consigli"
possono essere poco informative: nella frase "vorrei fondare una gilda sulla
musica", il topic utile per la ricerca è "musica". In parallelo, il sistema
prova a riconoscere vincoli espressi in linguaggio naturale, come una specifica
epoca, la modalità di partecipazione o la natura individuale o di gruppo della
Gilda. Dopo il recupero, i risultati vengono deduplicati e selezionati con
soglie dinamiche, evitando di passare al modello documenti troppo simili o poco
pertinenti.

La generazione usa Ollama come server locale. Il microservizio costruisce un
prompt aumentato con istruzioni, domanda utente e contesto recuperato, poi
invoca il modello e restituisce la risposta. Lo streaming NDJSON permette di
separare eventi di retrieval, token generati, warning ed errori, migliorando la
percezione di reattività nell'interfaccia. Il servizio include anche controlli
operativi, come stime di memoria disponibile, gestione dei timeout, rate
limiting e distinzione tra errori applicativi e dipendenze non raggiungibili.

Durante lo sviluppo è stato implementato anche un workflow agentico per guidare
la proposta di nuove Gilde. La tesi lo documenta come sperimentazione, non come
nucleo della soluzione finale. La ragione è pratica: a parità di hardware
locale, l'approccio agentico introduce più chiamate al modello, maggiore latenza
e maggiore complessità di controllo. Nel perimetro del prototipo, un RAG più
classico e deterministico risulta più stabile, più semplice da mantenere e più
coerente con i vincoli disponibili. La componente agentica rimane quindi una
possibile evoluzione, da valutare solo se produce un beneficio misurabile.

\section*{Valutazione dell'inferenza locale}

La seconda parte della tesi affronta la sostenibilità dell'inferenza locale.
Un sistema RAG aziendale può essere architetturalmente corretto, ma la sua
utilità dipende anche dal modello generativo usato nell'ultimo stadio della
pipeline. Il modello deve rientrare nella memoria disponibile, caricarsi in
tempi accettabili, iniziare a generare senza ritardi e mantenere una qualità
sufficiente per il caso d'uso. Per questo viene definito un protocollo
sperimentale dedicato ai modelli open-weight eseguiti localmente su Apple
Silicon.

L'esperimento usa un MacBook Pro con chip M1 Pro e 16 GB di memoria unificata.
Il runtime scelto è MLX, perché è progettato per Apple Silicon e consente di
convertire, quantizzare ed eseguire modelli linguistici in modo coerente con il
backend Metal. La matrice include modelli tra 3B e 14B parametri, divisi in
classi small, medium e large. Per i modelli piccoli vengono confrontate FP16,
Q8, Q6 e Q4; per quelli medi e grandi l'attenzione si concentra soprattutto
sulle varianti quantizzate, perché la quantizzazione diventa una condizione di
fattibilità entro il vincolo dei 16 GB.

Le metriche misurate sono dimensione su disco, memoria di picco, tempo di
caricamento, time-to-first-token, throughput di decodifica e throughput
end-to-end. A queste si aggiungono controlli automatici di qualità su subset
ridotti MMLU e GSM8K. La scelta di subset piccoli non mira a costruire una
classifica generale dei modelli; serve a individuare degradazioni macroscopiche,
problemi di formato e incompatibilità con il parser. I fallimenti non vengono
rimossi dall'analisi, perché in un sistema reale anche il formato dell'output è
parte del comportamento osservabile.

I risultati mostrano che, sui modelli piccoli, Q4 è il compromesso più
interessante. Riduce la dimensione degli artefatti di circa il 72\%, riduce il
picco di memoria di circa 67--70\% e aumenta il throughput di decodifica di
circa 2.6--2.8 volte rispetto a FP16. Nei controlli automatici ridotti non
emerge un degrado macroscopico. Q8 e Q6 sono varianti più conservative, ma nel
setup misurato non offrono un vantaggio qualitativo sufficiente a compensare il
maggiore uso di memoria.

Per i modelli medi e grandi la lettura cambia. Qwen3 8B e Mistral 7B diventano
utilizzabili in Q4 con throughput intorno a 32--35 token al secondo e picchi di
memoria inferiori a 5 GB. Gemma 3 12B e Qwen3 14B sono eseguibili in Q4, ma con
throughput intorno a 18 token al secondo e time-to-first-token superiore al
secondo. Queste configurazioni dimostrano la fattibilità tecnica, ma non sono
automaticamente preferibili in un'applicazione interattiva. Il margine residuo
di memoria deve infatti essere condiviso con Qdrant, FastAPI, Next.js, sistema
operativo e altri processi.

\section*{Risultati, limiti e sviluppi futuri}

La tesi conferma che un sistema RAG locale per un contesto aziendale è
realizzabile, ma richiede controllo simultaneo di più livelli. Non basta
scegliere un modello linguistico: occorre curare il corpus, progettare l'indice,
definire la logica di retrieval, controllare prompt e formato della risposta,
verificare la sostenibilità dell'inferenza e integrare il tutto con
l'applicazione esistente. Il contributo principale del lavoro è quindi un
prototipo integrato e valutato in condizioni realistiche, più che una singola
innovazione algoritmica.

Dal punto di vista architetturale, il prototipo mostra che i3Guild può essere
trasformato in una knowledge base interrogabile mantenendo PostgreSQL come
sorgente autoritativa e separando l'indice vettoriale dal database applicativo.
Dal punto di vista del retrieval, l'uso congiunto di embedding densi, segnali
sparsi, filtri strutturati, deduplicazione e selezione dinamica consente di
costruire un contesto più aderente al dominio rispetto a una ricerca puramente
semantica. Dal punto di vista sperimentale, la quantizzazione MLX rende
praticabili diversi modelli open-weight su hardware consumer Apple, con un
profilo particolarmente favorevole per i modelli piccoli in Q4.

Restano alcuni limiti. Il primo riguarda la valutazione applicativa del RAG: il
prototipo è stato implementato, ma manca una ground truth estesa costruita con
stakeholder aziendali per misurare groundedness, completezza, utilità e
allucinazioni sul corpus reale. Il secondo riguarda i dati: le Gilde sono
eterogenee per granularità, stile e completezza dei campi. Il terzo riguarda
l'esperimento MLX, eseguito su una sola macchina e con subset di qualità
ridotti; i risultati non possono essere generalizzati direttamente ad altri
hardware o runtime.

Gli sviluppi futuri seguono da questi limiti. La priorità è costruire un dataset
di domande realistiche sul dominio i3Guild, con risposte attese e documenti di
riferimento validati da persone interne all'azienda. Su questa base si potranno
confrontare versioni diverse della pipeline e misurare il contributo effettivo
di query rewriting, filtri, soglie e deduplicazione. Un secondo sviluppo
riguarda l'indicizzazione incrementale, che potrebbe sostituire la
reindicizzazione completa quando il corpus crescerà. Un terzo riguarda la
qualità del contesto recuperato, migliorabile con glossari aziendali, sinonimi
controllati, tag più espliciti e, in prospettiva, reranker.

La componente agentica potrà essere ripresa solo dopo il consolidamento del RAG
di base. Un agente potrebbe aiutare nella proposta di nuove Gilde, nella
raccolta di requisiti o nel collegamento con iniziative passate, ma la sua
adozione dovrebbe dipendere da criteri misurabili: riduzione del lavoro per
l'utente, qualità superiore delle bozze, controllo dei costi computazionali e
robustezza del formato prodotto. Nel perimetro attuale, la direzione più utile
non è aumentare la complessità dell'orchestrazione, ma rafforzare il legame tra
corpus, retrieval, valutazione e vincoli reali dell'infrastruttura locale.

\end{document}
